\documentclass[../../main.tex]{subfiles}

\begin{document}

\chapter{Pertukaran Kunci: Protokol Diffie-Hellman}

\begin{learningoutcome}
Setelah mempelajari bab ini, mahasiswa diharapkan mampu:
\begin{itemize}
    \item Menjelaskan persoalan distribusi kunci pada kriptografi simetri dan bagaimana protokol pertukaran kunci menyelesaikannya (Sub-CPMK 2.3).
    \item Menguraikan mekanisme kerja Protokol Pertukaran Kunci Diffie-Hellman (DH) (Sub-CPMK 2.3).
    \item Menerapkan perhitungan matematis Diffie-Hellman untuk menghasilkan kunci rahasia bersama (\textit{shared secret key}) (Sub-CPMK 2.3).
    \item Menganalisis landasan keamanan Diffie-Hellman berdasarkan sulitnya perhitungan logaritma diskrit (Sub-CPMK 2.3).
    \item Mengevaluasi kerentanan protokol Diffie-Hellman terhadap serangan aktif, khususnya \textit{Man-in-the-Middle Attack} (Sub-CPMK 2.3).
    \item Menjelaskan perluasan protokol Diffie-Hellman untuk komunikasi lebih dari dua pihak (Sub-CPMK 2.3).
\end{itemize}
\end{learningoutcome}

\subfile{section-01}
\subfile{section-02}
\subfile{section-03}
\subfile{section-04}
\section{Contoh Terhitung}

\begin{example}[Pertukaran kunci Diffie-Hellman dengan bilangan kecil]
Alice dan Bob menyepakati parameter publik $p = 11$ dan $g = 7$. Bilangan $7$ adalah akar primitif dari $11$ karena $7^1, 7^2, \dots, 7^{10} \pmod{11}$ menghasilkan semua nilai $1$ sampai $10$ (lihat Gambar~\ref{fig:contoh-numerik-dh}).
\begin{enumerate}
    \item Alice memilih kunci privat $a = 6$ dan menghitung kunci publiknya:
    \[ A = g^a \bmod p = 7^6 \bmod 11 = 4. \]
    \item Bob memilih kunci privat $b = 9$ dan menghitung kunci publiknya:
    \[ B = g^b \bmod p = 7^9 \bmod 11 = 8. \]
    \item Alice menerima $B = 8$ dan menghitung kunci rahasia bersama:
    \[ K = B^a \bmod p = 8^6 \bmod 11 = 3. \]
    \item Bob menerima $A = 4$ dan menghitung kunci rahasia bersama:
    \[ K = A^b \bmod p = 4^9 \bmod 11 = 3. \]
\end{enumerate}
Kedua pihak memperoleh kunci rahasia bersama $K = 3$ tanpa pernah mengirimkan nilai $K$ melalui saluran publik.
\end{example}

\begin{example}[Contoh dengan bilangan prima yang lebih besar]
Alice dan Bob menyepakati $p = 353$ dan $g = 3$; keduanya valid karena $3 < 353$ dan $3$ adalah akar primitif dari $353$.
\begin{enumerate}
    \item Alice memilih $a = 97$: \quad $A = g^a \bmod p = 3^{97} \bmod 353 = 40$.
    \item Bob memilih $b = 233$: \quad $B = g^b \bmod p = 3^{233} \bmod 353 = 248$.
    \item Alice menghitung: \quad $K = B^a \bmod p = 248^{97} \bmod 353 = 160$.
    \item Bob menghitung: \quad $K = A^b \bmod p = 40^{233} \bmod 353 = 160$.
\end{enumerate}
Penyadap yang melihat $p, g, A = 40$, dan $B = 248$ harus menyelesaikan persoalan logaritma diskrit untuk memperoleh $a$ atau $b$. Untuk $p$ sebesar ini perhitungan masih terjangkau komputer, tetapi untuk $p$ berukuran ratusan digit (misalnya 2048 bit) menjadi tidak praktis.
\end{example}


\begin{obeactivity}
\textbf{Aktivitas 11.1: Simulasi Perhitungan Diffie-Hellman}
1. Gunakan parameter publik $p = 353$ dan $g = 3$.
2. Alice memilih kunci privat $a = 97$, hitunglah nilai $A$.
3. Bob memilih kunci privat $b = 233$, hitunglah nilai $B$.
4. Hitunglah kunci rahasia bersama $K$ yang dihasilkan oleh Alice dan Bob.
5. Diskusikan bagaimana jika $p$ yang digunakan adalah bilangan kecil; apakah keamanan tetap terjamin?
\end{obeactivity}


\begin{obereflection}
\textbf{Latihan 11.1}
1. Mengapa pemilihan akar primitif $g$ sangat penting dalam protokol Diffie-Hellman?
2. Jelaskan perbedaan mendasar antara peran kunci publik dan kunci privat dalam RSA dibandingkan dengan dalam protokol Diffie-Hellman.
3. Bagaimana cara mencegah serangan \textit{Man-in-the-Middle} pada pertukaran kunci Diffie-Hellman? (Petunjuk: Pikirkan tentang tanda tangan digital atau sertifikat).
4. Sebuah organisasi memiliki 500 pegawai yang perlu berkomunikasi rahasia secara
   berpasangan memakai AES. Hitunglah berapa kunci rahasia yang harus dikelola bila
   setiap pasangan memakai kunci tersendiri. Bandingkan dengan jumlah kunci bila
   organisasi itu memakai pasangan kunci Diffie--Hellman sesaat untuk setiap sesi.
5. Jelaskan mengapa memperbesar modulus $p$ tidak menghalangi serangan
   \textit{Man-in-the-Middle}, padahal memperbesar $p$ menghalangi penyadapan pasif.
   Apa yang sesungguhnya menjadi celah keamanan pada serangan tersebut?
6. Pada protokol Diffie--Hellman sesaat (DHE), kunci privat $a$ dan $b$ dibuang
   setelah sesi berakhir. Jelaskan bagaimana kebiasaan itu melindungi rekaman
   percakapan yang sudah lama tersimpan pada saat kunci jangka panjang peladen
   akhirnya bocor.
\end{obereflection}

\begin{obereflection}
\textbf{Latihan 11.2}
1. Perhatikan $p = 23$, yang merupakan \textit{safe prime} karena $23 = 2 \times 11 + 1$
   dengan $q = 11$ juga prima. Tunjukkan bahwa $g = 2$ bukan akar primitif dari $23$,
   lalu hitunglah ordenya. Jelaskan mengapa $g$ semacam itu justru dianjurkan pada
   protokol yang memakai subgrup berorde prima.
2. Diketahui $p = 23$ dan $g = 5$. Tentukan seluruh akar primitif dari $23$ dan
   periksa bahwa jumlahnya sama dengan $\varphi(22)$. Tunjukkan bahwa $5$ termasuk
   di dalamnya.
3. Diketahui $p = 23$, $g = 5$, $a = 7$, dan $b = 15$. Hitunglah $A$, $B$, dan kunci
   bersama $K$ dari kedua sisi. Periksa hasil Anda dengan
   Contoh~\ref{ex:dh-kecil} dan dengan program pada Listing~\ref{lst:diffie-hellman}.
4. Modulus $p = 353$ memiliki $160$ akar primitif, yaitu sebanyak $\varphi(352)$.
   Jelaskan mengapa jumlah akar primitif modulo $p$ selalu sama dengan
   $\varphi(p-1)$, dan apa akibat praktisnya bagi pihak yang memilih $g$ secara acak.
5. Dengan algoritma langkah raksasa-bayi, penyerang memerlukan sekitar
   $2\sqrt{p}$ langkah. Untuk modulus 2048 bit, tunjukkan bahwa $\sqrt{2^{2048}}$
   berukuran sekitar $1{,}8 \times 10^{308}$. Bila satu komputer dapat melakukan
   $10^{9}$ perkalian per detik, perkirakan berapa tahun waktu yang diperlukan.
   Bandingkan dengan umur alam semesta yang sekitar $1{,}4 \times 10^{10}$ tahun.
6. Pada contoh $p = 353$, penyerang menemukan $a = 97$ hanya dengan $24$
   perkalian langkah raksasa-bayi, jauh lebih sedikit daripada $97$ langkah yang
   diperlukan pencarian menyeluruh. Bila Alice memilih kunci privat yang lebih
   besar, misalnya $a = 300$, jumlah perkaliannya memang bertambah, tetapi tetap
   berorde $\sqrt{p}$ dan bukan berorde $a$. Jelaskan mengapa demikian, dan apa
   yang sebenarnya diukur oleh perkiraan $2\sqrt{p}$ itu.
7. Rancanglah satu skema sederhana yang memungkinkan Alice dan Bob memastikan
   keaslian nilai pertukaran mereka dengan memakai tanda tangan digital. Jelaskan
   bagian mana dari skema Anda yang membuat Mallory tidak dapat menggantikan
   $g^{a}$ dengan $g^{m}$.
\end{obereflection}

\section{Rangkuman}
\begin{itemize}
    \item Persoalan distribusi kunci pada kriptografi simetris muncul karena $n$
          pengguna memerlukan $\binom{n}{2} = n(n-1)/2$ kunci rahasia berbeda;
          untuk 1.000 pengguna jumlahnya sudah mencapai 499.500 kunci. Persoalan
          itu diselesaikan dengan kriptografi hibrida atau dengan protokol
          pertukaran kunci seperti Diffie--Hellman (1976).
    \item Kriptografi hibrida mengenkripsi kunci sesi yang pendek dengan RSA,
          lalu memakai kunci sesi itu untuk memproses seluruh isi pesan dengan
          AES. Enam langkahnya dirumuskan sebagai
          $C_K = E_{\text{RSA},\text{PK}_{\text{Bob}}}(K)$ dan
          $C_M = E_{\text{AES},K}(M)$, yang dipulihkan Bob dengan
          $K = D_{\text{RSA},\text{SK}_{\text{Bob}}}(C_K)$ dan
          $M = D_{\text{AES},K}(C_M)$.
    \item Diffie--Hellman memungkinkan dua pihak menyepakati kunci rahasia bersama
          melalui saluran publik tanpa pernah mengirimkan kunci itu sendiri.
    \item Parameter publiknya adalah bilangan prima besar $p$ dan akar primitif
          $g$ dari $p$. Akar primitif berarti orde $g$ modulo $p$ sama dengan
          $p-1$; bila syarat itu tidak dipenuhi, himpunan kunci publik yang
          mungkin menjadi kecil dan dapat didaftar seluruhnya. Untuk $g = 3$
          modulo 11, hanya lima kunci publik yang mungkin.
    \item Prosedurnya: Alice menghitung $A = g^a \bmod p$, Bob menghitung
          $B = g^b \bmod p$, lalu keduanya memperoleh kunci yang sama
          $K = B^a = A^b = g^{ab} \bmod p$. Kesamaan itu berlaku karena
          $(g^b)^a = (g^a)^b$; yang menyamakan kedua kunci justru urutan langkah
          yang berbeda di kedua sisi.
    \item Contoh: $p = 353$, $g = 3$, $a = 97$, dan $b = 233$ menghasilkan
          $A = 40$, $B = 248$, serta kunci bersama $K = 160$. Contoh dari
          perkakas daring, $G = 1601$ dan $N = 4789$, menghasilkan $A = 2716$,
          $B = 1302$, dan $K = 4257$.
    \item Keamanannya bersandar pada sulitnya menghitung logaritma diskret:
          penyadap yang mengetahui $p, g, A, B$ tetap sukar menemukan $a$ atau
          $b$ bila $p$ sangat besar. Algoritma langkah raksasa-bayi menurunkan
          biaya dari $p-1$ menjadi sekitar $2\sqrt{p}$ langkah; untuk modulus
          2048 bit $\sqrt{p}$ berukuran sekitar $2^{1024}$, sehingga panjang
          minimum yang diterima hari ini adalah 2048 bit.
    \item Diffie--Hellman aman terhadap penyadapan pasif tetapi rentan terhadap
          serangan aktif \textit{Man-in-the-Middle}, sebab tidak ada satu pun
          langkahnya yang memeriksa identitas lawan bicara. Mallory menyepakati
          $K_1 = g^{am} \bmod p$ dengan Alice dan $K_2 = g^{bm} \bmod p$ dengan
          Bob, tanpa perlu mengetahui $a$ maupun $b$. Penangkalnya adalah
          otentikasi, misalnya melalui tanda tangan digital atau sertifikat.
    \item Pemilihan parameter menentukan keamanan: $p$ harus cukup panjang dan
          idealnya berbentuk \textit{safe prime} $p = 2q+1$, $g$ harus
          membangkitkan subgrup berorde prima besar, dan $a$, $b$ harus
          dibangkitkan dengan keacakan yang benar. Serangan Logjam (2015)
          memperlihatkan bahwa penurunan mutu parameter dapat menjatuhkan
          protokol tanpa menyentuh matematikanya.
    \item Protokol ini dapat diperluas ke tiga pihak dengan pola melingkar,
          menghasilkan $K = g^{abc} \bmod p$, tetapi jumlah pengirimannya tumbuh
          secepat $n(n-1)$ sehingga hanya praktis untuk kelompok kecil.
    \item Diffie--Hellman sesaat (DHE) dan versi kurva eliptiknya (ECDH) memberi
          \textit{forward secrecy}: kebocoran kunci jangka panjang tidak membuka
          rekaman percakapan yang sudah terjadi. Karena alasan itu TLS 1.3
          mewajibkan pertukaran kunci sesaat, dan ECDH dengan kunci 256 bit
          menggantikan Diffie--Hellman modulus 3072 bit.
    \item Protokol ini dipakai pada TLS/HTTPS, SSH, IPsec, dan aplikasi pesan
          modern. Polanya selalu sama: Diffie--Hellman menetapkan kunci,
          sedangkan algoritma kunci-simetri memproses isi pesan.
\end{itemize}


\begin{obeassessment}
\textbf{Kuis Bab 11}
1. Jelaskan langkah-langkah matematis yang terjadi dalam pertukaran kunci Diffie-Hellman.
2. Apa yang dimaksud dengan masalah logaritma diskrit dan mengapa hal ini menjadi fondasi keamanan DH?
3. Gambarkan alur serangan \textit{Man-in-the-Middle} pada protokol DH dan jelaskan mengapa hal tersebut bisa terjadi.
\end{obeassessment}
\begin{competencychecklist}
\begin{tabularx}{\linewidth}{|X|c|}
\hline
Indikator Kompetensi & Tercapai \\
\hline
Saya dapat menjelaskan masalah distribusi kunci pada simetri & [ ] \\
\hline
Saya dapat menguraikan prosedur pertukaran kunci Diffie-Hellman & [ ] \\
\hline
Saya dapat menghitung kunci rahasia bersama $K$ menggunakan parameter DH & [ ] \\
\hline
Saya dapat menganalisis risiko serangan MITM pada protokol DH & [ ] \\
\hline
Saya dapat menjelaskan implementasi DH pada protokol TLS/SSL & [ ] \\
\hline
\end{tabularx}
\end{competencychecklist}


\end{document}
